\section{Methods}\label{sec:methods}
% Budget: 5-5.5 pages. Generated from docs/decisions.md (sections 2-5).

Let $y_t\in\mathbb{R}^N$ be the vector of weekly returns ($N=18$). All methods are estimated on rolling windows of $w=104$ weeks; each estimate is dated at the last week of its window. For the comparison, all methods use the same 179 windows (step of four weeks, first window ending 1 March 2013).

\subsection{Diebold--Yilmaz connectedness}
We fit a VAR($p$), $y_t = c + \sum_{k=1}^p A_k y_{t-k} + \varepsilon_t$, $\mathbb{E}[\varepsilon_t\varepsilon_t'] = \Sigma$, by OLS, with $p\in\{1,2\}$ selected by BIC in each window. The cap $p\le 2$ is an adaptation: with $N=18$ and about 100 observations per window, $p=3$ would mean 55 regressors per equation. With MA coefficients $\Psi_h$ ($\Psi_0=I$), the $H$-step generalized variance share \citep{pesaran1998generalized} is
\begin{equation}
\theta_{ij}(H) = \frac{\sigma_{jj}^{-1}\sum_{h=0}^{H-1}(e_i'\Psi_h\Sigma e_j)^2}{\sum_{h=0}^{H-1} e_i'\Psi_h\Sigma\Psi_h' e_i},
\qquad \tilde\theta_{ij} = \theta_{ij}\Big/\textstyle\sum_{j}\theta_{ij},
\end{equation}
which is invariant to the ordering of the variables. Following \citet{diebold2014network}, $C^{\text{from}}_i = 100\sum_{j\ne i}\tilde\theta_{ij}$, $C^{\text{to}}_j = 100\sum_{i\ne j}\tilde\theta_{ij}$, net connectedness $C^{\text{net}}_i = C^{\text{to}}_i - C^{\text{from}}_i$, pairwise net $100(\tilde\theta_{ji}-\tilde\theta_{ij})$, and total connectedness $C = \frac{100}{N}\sum_{i\ne j}\tilde\theta_{ij}$. The baseline horizon is $H=4$ weeks. The base paper uses daily data, 100-day windows and $H=12$ days; we keep the logic (a horizon of about one month, a window long enough for a stable VAR) on the weekly scale.

\subsection{Sparse VAR (lasso)}
As a bridge between the dense DY VAR and the sparse network estimators, we estimate each VAR equation by the lasso \citep{tibshirani1996lasso} on lags $1,\dots,4$ (standardized regressors, unpenalized intercept), so that the penalty performs lag selection. The penalty is chosen per equation by time-series cross-validation (five expanding folds, test blocks always after the training block; random $K$-fold would leak future information), over 40 log-spaced values; the extended BIC \citep{chen2008ebic} is an alternative. The same GFEVD and connectedness measures are computed on the lasso VAR, so differences with DY isolate the effect of regularization.

\subsection{NETS}
NETS \citep{barigozzi2019nets} jointly estimates a sparse VAR(1) and a sparse concentration matrix $K=\Sigma^{-1}$ of the innovations by adaptive lasso \citep{zou2006adaptive, peng2009space}. Since the \texttt{nets} package requires the penalty to be supplied, we select it by BIC over 15 log-spaced values. The Granger network is given by the non-zero entries of $A_1$, the contemporaneous network by the partial correlations $-K_{ij}/\sqrt{K_{ii}K_{jj}}$, and the long-run network by the partial correlations of $K_{LR} = (I-A(1))'K(I-A(1))$. Each window is standardized before estimation.

\subsection{FNETS}
FNETS \citep{barigozzi2024fnets, owens2024fnetsR} assumes $y_t = \chi_t + \xi_t$, where the common component $\chi_t$ is driven by $q$ dynamic factors and the idiosyncratic component $\xi_t$ follows a sparse VAR. We use the unrestricted (generalized dynamic) factor model, select $q$ with the \citet{hallin2007determining} information criterion (criterion IC5, the package default), estimate the idiosyncratic VAR(1) by lasso with the penalty chosen by the package's BIC, and estimate contemporaneous and long-run partial correlations of $\xi_t$ with the package's CLIME-type estimator \citep{cai2016estimating} and adaptive thresholding. The share of variance explained by the common component is $\mathrm{var}(\chi_{it})/\mathrm{var}(y_{it})$. Applying the GFEVD to the idiosyncratic VAR yields \emph{idiosyncratic connectedness}; comparing it with DY measures how much connectedness is due to the common factor. With $N=18$ and $T=104$ the estimated idiosyncratic innovation covariance is often not positive definite; we project it onto the positive-definite cone by clipping eigenvalues at 1\% of the mean absolute eigenvalue, and clip partial correlations to $[-1,1]$ (the support of the network is unchanged).

\subsection{Comparable measures and comparison criteria}
Every method is mapped to the same objects: a VAR on which the same GFEVD connectedness is computed; a directed Granger network (sum over lags of the absolute coefficients); and contemporaneous and long-run partial-correlation networks (for DY and the lasso VAR, implied by the estimated VAR). We compare methods on three criteria:
\begin{enumerate}
\item \emph{Density}: share of non-zero links in each network; for the dense GFEVD table, share of pairwise shares above $1/N$; and total connectedness.
\item \emph{Agreement}: correlation of total and net connectedness series across methods; per-window Spearman correlation between the off-diagonal entries of two methods' matrices; Jaccard similarity of their strongest 10\% of links.
\item \emph{Stability}: variability of total connectedness, and Spearman correlation and top-link Jaccard similarity between consecutive windows' matrices.
\end{enumerate}
Finally, the share of DY connectedness that disappears once the common factor is removed is $1 - C^{\text{FNETS}}/C^{\text{DY}}$, averaged over windows.

\subsection{Robustness}
We rerun all methods with windows of 78 and 130 weeks, horizons of 2 and 8 weeks, without SPY, and excluding 2020. For the last, every rolling window containing any week of 2020 is dropped (the resulting gap is not interpolated), and static full-period connectedness is estimated separately for 2011--2019 and 2021--2026.
